\subsection{An explicit Hecke calculation on the compact unitary space}
\label{sec:u3-hecke}

Write $E=\mathbf Q(\omega)$, where $\omega^2+\omega+1=0$, and
$O=\mathbf Z[\omega]$.  Let $G=U(h)$ for
$h(z,w)=\sum_{i=1}^3\overline{z_i}w_i$, and let $K$ be the full
finite adelic stabilizer of $O^3$.  All local Hecke measures below give
the corresponding maximal compact subgroup volume one.  The following
calculation concerns this precise group, level, and coefficient space;
the passage from other automorphic representations to this space is the
separate coverage theorem.

\begin{theorem}[K17-U3-HECKE]\label{thm:u3-hecke}
Let $V$ be the irreducible compact $U(3)$ representation of highest
weight $(5,0,-5)$.  The algebraic modular-form space $M(G,K,V)$ has
dimension one.  For $p=7,13$, choose the split factor determined by
\[
 \pi_7=3+\omega,\qquad \pi_{13}=4+\omega,
 \qquad N(\pi_p)=p.
\]
Let $T_p^+$ be the unnormalized positive minuscule Hecke operator,
corresponding at that factor to
$\mathrm{GL}_3(\mathbf Z_p)\operatorname{diag}(p,1,1)
 \mathrm{GL}_3(\mathbf Z_p)$, and write $t_p$ for the normalized
Satake trace.  Then
\[
 T_p^+=p\,t_p,\qquad A_p:=p^5T_p^+=p^6t_p,
 \qquad
 \begin{array}{c|r|r}
 p&\#\text{Hecke cosets}&A_p\\ \hline
 7&57&157899\\
 13&183&6110859
 \end{array}
 \tag{UH.1}
\]
In particular, the characteristic polynomials of the scaled operators
are $X-157899$ and $X-6110859$.  Their values at the comparison
integers satisfy
\begin{align*}
 7^{12}+7^6+1-157899
   &=13841246952\\
   &=2^3 3^5\cdot13\cdot43\cdot47\cdot271,\\
 13^{12}+13^6+1-6110859
   &=23298083838432\\
   &=2^5 3^5\cdot13\cdot47\cdot337\cdot14551.
\end{align*}
\end{theorem}

\begin{proof}
We give the class-number argument, the coefficient model, and a
complete finite calculation.

\paragraph{The lattice class number.}
The Eisenstein lattice in $\mathbf C$ has covering-radius squared
$1/3$: its hexagonal Voronoi cell has vertices at that distance from
the origin.  Hence every complex number is congruent modulo $O$ to
one of squared absolute value at most $1/3$.  In particular $O$ is
norm-Euclidean, and therefore a principal ideal domain.

Let $L$ be any positive integral unimodular Hermitian $O$-lattice of
rank three.  Choose a shortest primitive vector, then a shortest
primitive vector in its orthogonal projection, and then a basis of the
remaining rank-one projection.  Primitivity and the PID property let
these vectors be lifted to a basis of $L$.  Write $a,b,c$ for the
successive Gram--Schmidt squared lengths.  Reducing the coefficient
of the second vector along the first modulo $O$ produces a vector of
norm at most $b+a/3$.  Minimality of the first vector gives
$b\geq 2a/3$.  The same argument in the projected rank-two lattice
gives $c\geq 2b/3$.  Consequently
\[
 1=\det L=abc\geq \frac{8a^3}{27}.
 \tag{UH.2}
\]
The minimum $a$ is a positive integer, so $a\leq3/2$ forces $a=1$.
A norm-one vector splits off integrally and orthogonally: subtract
its integral Hermitian pairing from any vector.  Its complement is
again integral unimodular.  In rank two, the same reduction gives
$1\geq2a^2/3$, again forcing its integral minimum to be one.  The
remaining rank-one form also has norm one.  Thus $L\simeq O^3$.

The argument proves the analogous rank-two assertion as well.  Every
adelic translate of $O^3$ intersects $E^3$ in an integral unimodular
global lattice; conversely its localizations give the translated
lattices.  The preceding isometry therefore proves
\[
 G(\mathbf Q)\backslash G(\mathbf A_f)/K=\{1\}.
 \tag{UH.3}
\]
The integral stabilizer is
$\Gamma=\mu_6^3\rtimes S_3$: a norm-one integral vector has precisely
one unit coordinate, since its norm is a sum of three nonnegative
integral norms; orthogonality then forces a unitary integral matrix
to be monomial.

\paragraph{The invariant harmonic polynomial.}
Realize $V$ as the harmonic polynomials of bidegree $(5,5)$ in
$z_1,z_2,z_3$ and their conjugates.  Here the action is
\[
       (\rho(g)f)(z)=f(g^{-1}z),\qquad
       \Delta=\sum_i\partial_{z_i}\partial_{\overline z_i}.
 \tag{UH.4}
\]
For clarity, this realization can be checked directly.  Contraction
from bidegree $(5,5)$ onto $(4,4)$ is surjective.  If $R=\sum z_i\bar z_i$
and $F$ has bidegree $(4,4)$, a preimage is
$\sum_{j=0}^4 c_jR^{j+1}\Delta^jF$, where
$c_0=1/11$ and $c_j=-c_{j-1}/((j+1)(11-j))$; the product rule verifies
the assertion.  Its kernel therefore has dimension $21^2-15^2=216$.
It contains the highest-weight vector $\bar z_1^5z_3^5$ of weight
$(5,0,-5)$, whose irreducible representation has Weyl dimension
$6\cdot6\cdot12/2=216$.  This proves the claimed realization.  The
dual has the same highest weight, so using the dual coefficient
convention makes no difference here.

Invariance under the three independent sixth roots of unity forces
the exponents of $z_i$ and $\bar z_i$ to agree: their difference is
divisible by six and lies between $-5$ and $5$.  Put $x_i=z_i\bar z_i$.
The $\Gamma$-invariant polynomials are exactly the symmetric
homogeneous degree-five polynomials in the $x_i$.  If $m_\lambda$
denotes the sum over the distinct permutations of the monomial with
exponent partition $\lambda$, the Laplacian is
\[
 \sum_i(x_i\partial_{x_i}^2+\partial_{x_i}).
\]
Its matrix, with source partitions $(5)$, $(4,1)$, $(3,2)$,
$(3,1,1)$, $(2,2,1)$ and target partitions $(4)$, $(3,1)$,
$(2,2)$, $(2,1,1)$, is
\[
 \begin{pmatrix}
 25&2&0&0&0\\
 0&16&4&1&0\\
 0&0&18&0&1\\
 0&0&0&9&8
 \end{pmatrix}.
 \tag{UH.5}
\]
It has rank four.  Thus $V^\Gamma$ is the line spanned by
\[
 f=2m_{(5)}-25m_{(4,1)}+20m_{(3,2)}
       +320m_{(3,1,1)}-360m_{(2,2,1)},
 \qquad f(1,0,0)=2.
 \tag{UH.6}
\]
An algebraic modular form satisfies
$F(\gamma gk)=\rho(\gamma)F(g)$; evaluation at the unique adelic
class identifies $M(G,K,V)$ with $V^\Gamma$.  This proves its dimension.

\paragraph{The neighbor correspondence and its direction.}
For a split prime $p=\pi\bar\pi$, suppose an integral matrix $U$
satisfies
\[
 U^*U=pI,\qquad \operatorname{rank}(U\bmod\pi)=1.
 \tag{UH.7}
\]
Then $g=U/\pi$ is unitary.  Both $U$ and $U^*$ are integral and
$U^{-1}=U^*/p$, so its $\pi$-elementary divisors are zero or one;
rank one modulo $\pi$ makes them $(0,1,1)$.  Thus $g$ has elementary
divisors $(-1,0,0)$ at $\pi$, $(0,0,1)$ at $\bar\pi$, and is integral
invertible elsewhere.  It initially represents the \emph{negative}
minuscule correspondence at the chosen factor.  Its right coset is
determined by
\[
    \ell=\operatorname{im}(U\bmod\pi)\in\mathbf P^2(\mathbf F_p).
 \tag{UH.8}
\]
Conversely all these lines classify the right cosets: they specify
the one-dimensional enlargement of $O_\pi^3$ inside
$\pi^{-1}O_\pi^3$, and Hermitian duality specifies the lattice at
$\bar\pi$.  There are $p^2+p+1$ such cosets.

For each representative $g_\ell$, the contribution to the Hecke action
on $f$, evaluated at $e_1$, is
\[
 (\rho(g_\ell)f)(e_1)
   =p^{-5}f\bigl(N(U_{11}),N(U_{12}),N(U_{13})\bigr).
 \tag{UH.9}
\]
These are \emph{first-row} norms, because the action uses $g^{-1}$.
Changing the representative on the right by an element of $\Gamma$
does not change this expression.  The subgroup of $\Gamma$ preserving
$\mathbf Ce_1$ acts on the lines in (UH.8), and the expression is
constant on each of its orbits.

The $p+1$ lines with first coordinate zero can be represented inside
the rank-two unitary group on the last two coordinates.  The proved
rank-two class-number assertion supplies those rational representatives.
They fix $e_1$.  The line $\mathbf F_pe_1$ has representative
$\operatorname{diag}(\bar\pi/\pi,1,1)$, which acts on $e_1$ by a scalar
of complex absolute value one.  These $p+2$ terms all have the value
given by the norm pattern $(p,0,0)$ in (UH.9).

For the other lines, normalize their first coordinate to one.  At
$p=7$, reduction maps $\mu_6$ onto $\mathbf F_7^\times$.  The remaining
two orbits have representatives $(1,1,0)$ and $(1,1,1)$ and sizes
$12$ and $36$.  At $p=13$, the image of $\mu_6$ is precisely the six
nonzero squares.  The five remaining orbits have representatives
\[
 (1,1,0),\ (1,2,0),\ (1,1,1),\ (1,1,2),\ (1,2,2)
\]
and sizes $12,12,36,72,36$.  The element $2$ is a nonsquare modulo
$13$.  Independent unit multiplications of the last two coordinates
and their interchange prove these exhaustive orbit descriptions.

\paragraph{Seven matrices give the complete calculation.}
Here are matrices satisfying (UH.7) for these representatives in the
order just listed.  Their reductions have exactly the indicated image
lines.  Every identity is checked directly using $\omega^2=-\omega-1$.
For $p=7$ they are
\[
 \begin{pmatrix}
 -2-2\omega&-2-\omega&0\\
 1-\omega&2\omega&0\\
 0&0&-3-\omega
 \end{pmatrix},\qquad
 \begin{pmatrix}
 -2-\omega&-2-\omega&-1-\omega\\
 -2-\omega&1&1+2\omega\\
 1&-2-\omega&1+2\omega
 \end{pmatrix}.
 \tag{UH.10}
\]
For $p=13$ the two representatives with one zero coordinate are
\[
 \begin{pmatrix}
 -4-2\omega&-1-\omega&0\\
 -\omega&-2+2\omega&0\\
 0&0&-4-\omega
 \end{pmatrix},\qquad
 \begin{pmatrix}
 -3-3\omega&-2-2\omega&0\\
 2\omega&-3\omega&0\\
 0&0&-4-\omega
 \end{pmatrix}.
 \tag{UH.11}
\]
The three representatives with all coordinates nonzero are
\begin{gather*}
 \begin{pmatrix}
 -3-2\omega&-2-\omega&-2-\omega\\
 1-\omega&-2-\omega&1+3\omega\\
 1-\omega&1+3\omega&-2-\omega
 \end{pmatrix},\qquad
 \begin{pmatrix}
 -3-3\omega&-2-\omega&-1-\omega\\
 \omega&-2-\omega&3\\
 1-\omega&-1+2\omega&1+2\omega
 \end{pmatrix},\\[3pt]
 \begin{pmatrix}
 -3-\omega&-2-\omega&-2-\omega\\
 -2-\omega&-\omega&3+3\omega\\
 -2-\omega&3+3\omega&-\omega
 \end{pmatrix}.
 \tag{UH.12}
\end{gather*}
Together with the $p+2$ elementary terms, these matrices and their
unit/permutation translates account for every Hecke coset.  Substitution
in (UH.6) gives the following complete table; equal norm patterns have
been combined.
\[
\begin{array}{c|c|r|r}
 p&\text{sorted first-row norms}&\text{multiplicity}&f\text{-value}\\\hline
 7&(7,0,0)&9&33614\\
 7&(4,3,0)&12&-4606\\
 7&(3,3,1)&36&1904\\\hline
 13&(13,0,0)&15&742586\\
 13&(12,1,0)&12&16406\\
 13&(9,4,0)&12&-256594\\
 13&(9,3,1)&72&40976\\
 13&(7,3,3)&72&14096
\end{array}
 \tag{UH.13}
\]
Dividing each weighted sum by $f(e_1)=2$ gives $157899$ and
$6110859$.  By (UH.9), these are $p^5$ times the eigenvalues of the
negative minuscule Hecke operator.

\paragraph{The positive operator and Satake factor.}
Complex conjugation sends the unitary matrices $g_\ell=U_\ell/\pi$
to representatives of the positive minuscule correspondence at the
chosen split factor.
It preserves $f$ and all first-row norms.
Thus the positive and negative operators have the same eigenvalue
on this invariant line.  This explicitly resolves their direction;
it is not inferred merely from the two operators having the same
number of cosets.

With maximal compact volume one, the Satake transform of the positive
$\mathrm{GL}_n$ minuscule characteristic function is
$p^{(n-1)/2}$ times the standard trace
\cite[\S3, (3.13)--(3.14), author preprint p.~8]{K17GrossSatake1998}.
For $n=3$ it is $p\,t_p$.  Hence (UH.13) gives
$A_p=p^5T_p^+=p^6t_p$, with no stabilizer-volume factor.  Expanding
the two comparison values and multiplying the displayed prime factors
proves the remaining assertions.
\end{proof}

\paragraph{Exact reproducibility.}
The proof above reduces the numerical result to seven small matrices,
the Laplacian matrix, and the eight rows of (UH.13).  The standalone
\path{work/hecke_level_one/u3_orbit_certificate.py} contains these
seven matrices, verifies their exact Gram and reduction identities,
enumerates the disjoint unit orbits, and reproduces the table and
prime factorizations without reading any input file.  The complete
neighbor enumeration is also independently replayed in
\path{work/hecke_level_one/tripod_independent_check.py}.
This standard-library verifier imports neither discovery program nor
a symbolic algebra package.  It reads the complete frozen neighbor
data in \texttt{u3\_weight5\_exact\_exploration.json}, checks $U^*U=pI$
and every distinct line of $\mathbf P^2(\mathbf F_p)$, verifies the
harmonic equation and a nonzero rank-four minor, and computes the
eigenvalues by exact integers.  It also checks the eigenvector identity
at five additional integral vectors for each prime.  The associated
JSON report records $57$ and $183$ checked cosets and the values in
(UH.1).  The original enumeration in \texttt{compute\_u3\_weight5.py}
gives the same result.  These are finite exact computations, not
machine-checked proofs in Lean; the complete mathematical correspondence
and normalization are the ones proved above.

No assertion about coverage of a different level, a different inner
form, or all conjugate-self-dual cusp forms is included in this theorem.
Those hypotheses must be supplied before using these integers in an
arithmetic divisibility criterion.
